\documentclass[14pt]{extarticle}
\usepackage{fontspec}
\usepackage{polyglossia}
\usepackage{amsmath}
\usepackage{amssymb}
\usepackage{xcolor}
\usepackage{fancyhdr}
\usepackage{graphicx}
\usepackage{listings}
\usepackage[most]{tcolorbox}
% \usepackage{setspace}
\usepackage{pifont}
\usepackage{enumitem}
\usepackage{titlesec}
\usepackage[bottom]{footmisc}
\usepackage{titling}
\usepackage{minted}
\usepackage{etoolbox}
\usepackage{array}
\usepackage{longtable} % page-breakable tables (must load before polyglossia/bidi RTL setup)
\usepackage{extsizes}
\usepackage{float}
\usepackage{booktabs}
\usepackage{hyperref}
\usepackage[onehalfspacing]{setspace}
\newif\ifwhiteblack
\ifdefined\setwhiteblack
  \setwhiteblack
\fi

\ifwhiteblack
  \usepackage[a4paper, landscape, top=1.25cm, bottom=1.25cm, left=1.25cm, right=1.25cm, headheight=1.5cm, headsep=0.5cm]{geometry}
\else
  \usepackage[a4paper, top=2.7cm, bottom=2.7cm, left=2.5cm, right=2.5cm, headheight=1.5cm, headsep=0.5cm]{geometry}
\fi
\usepackage{tikz}
\usepackage{datetime2}

\usetikzlibrary{automata, positioning, arrows.meta, arrows}
\providecommand{\LTcaptype}{table}

\newfontfamily\emoji{Segoe UI Emoji}

\pagestyle{fancy}

\setmainlanguage[numerals=western]{arabic}
\setotherlanguage{english}
\newfontfamily\arabicfont[Script=Arabic]{Calibri}
\newfontfamily\arabicfonttt[Script=Arabic]{Courier New}



\ifwhiteblack
\lstset{
  language=[Sharp]C,
  numbers=left,
  stepnumber=1,
  numbersep=8pt,
  frame=single,
  basicstyle=\ttfamily\small,
  keywordstyle=\bfseries,
  stringstyle=\ttfamily,
  commentstyle=\itshape
}
\else
\lstset{
  language=[Sharp]C,
  numbers=left,
  stepnumber=1,
  numbersep=8pt,
  frame=single,
  basicstyle=\ttfamily\small,
  keywordstyle=\color{blue},
  stringstyle=\color{red},
  commentstyle=\color{green!50!black}
}
\fi

\newif\ifdetailed
\ifdefined\setdetailed
  \setdetailed
\fi

\newif\ifwithsols
\ifdefined\setwithsols
  \setwithsols
\fi

\ifwhiteblack
  \AtBeginDocument{
    \LARGE
    \usemintedstyle{bw}
  }
\fi

% Default minted look: compact code that wraps long lines.
% Not applied in black-and-white landscape mode, which already uses \LARGE.
\ifwhiteblack\else
  \setminted{fontsize=\small, breaklines=true}
\fi

% unified tcolorboxes for articles
\ifwhiteblack
\tcbset{colback=white, colframe=black, fonttitle=\bfseries, boxrule=1pt}
\newtcolorbox{boxDef}[1][]{title={{\emoji📘} تعريف\ifx\\#1\\\else ~#1\fi :}}
\newtcolorbox{boxExercise}[1][]{title={{\emoji🧩} تمرين\ifx\\#1\\\else ~#1\fi :}}
\newtcolorbox{boxExample}[1][]{title={{\emoji📝} مثال\ifx\\#1\\\else ~#1\fi :}}
\newtcolorbox{boxNote}[1][]{title={{\emoji✨} ملاحظة\ifx\\#1\\\else ~#1\fi :}}
\newtcolorbox{boxAttention}[1][]{title={{\emoji🔔} تنبيه\ifx\\#1\\\else ~#1\fi :}}
\newtcolorbox{boxWarning}[1][]{title={{\emoji⚡} ملاحظة هامة\ifx\\#1\\\else ~#1\fi :}}
\newtcolorbox{boxSolution}[1][]{title={{\emoji✅} حل\ifx\\#1\\\else ~#1\fi :}}
\newtcolorbox{boxSymbol}[1][]{title={{\emoji🔣} رمز\ifx\\#1\\\else ~#1\fi :}}
\newtcolorbox{boxHint}[1][]{title={{\emoji💡} تلميح\ifx\\#1\\\else ~#1\fi :}}
\else
\tcbset{colback=white, colframe=black, fonttitle=\bfseries, boxrule=0.8pt}
\newtcolorbox{boxDef}[1][]{colback=blue!5!white,colframe=blue!75!black,
  title={{\emoji📘} تعريف\ifx\\#1\\\else ~#1\fi :}}
\newtcolorbox{boxExercise}[1][]{colback=cyan!5!white,colframe=cyan!70!black,
  title={{\emoji🧩} تمرين\ifx\\#1\\\else ~#1\fi :}}
\newtcolorbox{boxExample}[1][]{colback=yellow!5!white,colframe=orange!90!black,
  title={{\emoji📝} مثال\ifx\\#1\\\else ~#1\fi :}}
\newtcolorbox{boxNote}[1][]{colback=gray!10!white,colframe=black,
  title={{\emoji✨} ملاحظة\ifx\\#1\\\else ~#1\fi :}}
\newtcolorbox{boxAttention}[1][]{colback=magenta!10!white,colframe=magenta!80!black,
  title={{\emoji🔔} تنبيه\ifx\\#1\\\else ~#1\fi :}}
\newtcolorbox{boxWarning}[1][]{colback=red!5!white,colframe=red!75!black,
  title={{\emoji⚡} ملاحظة هامة\ifx\\#1\\\else ~#1\fi :}}
\newtcolorbox{boxSolution}[1][]{colback=green!5!white,colframe=green!60!black,
  title={{\emoji✅} حل\ifx\\#1\\\else ~#1\fi :}}
\newtcolorbox{boxSymbol}[1][]{colback=purple!5!white,colframe=purple!70!black,
  title={{\emoji🔣} رمز\ifx\\#1\\\else ~#1\fi :}}
\newtcolorbox{boxHint}[1][]{colback=teal!5!white,colframe=teal!60!black,
  title={{\emoji💡} تلميح\ifx\\#1\\\else ~#1\fi :}}
\fi


\ifwhiteblack
\tcbset{simplecode/.style={ colback=white, colframe=black, boxrule=1pt, arc=2pt, left=4pt,right=4pt,top=4pt,bottom=4pt}}
\else
\tcbset{simplecode/.style={ colback=gray!5, colframe=black!50, boxrule=0.4pt, arc=2pt, left=4pt,right=4pt,top=4pt,bottom=4pt}}
\fi
% Compact code boxes: tighter line spacing (the preamble uses onehalfspacing) and small vertical gaps.
% Must come AFTER the simplecode definition above.
\tcbset{simplecode/.append style={fontupper=\setstretch{0.95}, top=4pt, bottom=4pt, before skip=6pt, after skip=6pt}}
\newenvironment{boxCode}{\begin{tcolorbox}[simplecode]}{\end{tcolorbox}}

% Code cell used inside tables (keeps code left-to-right inside RTL tables).
% Use inside a fixed-width p{...} column: \begin{pycell}\begin{minted}...\end{minted}\end{pycell}
\newtcolorbox{pycell}{frame empty,colback=white,boxsep=0pt,left=2pt,right=2pt,top=2pt,bottom=2pt,width=\linewidth}

\newcolumntype{C}[1]{>{\centering\arraybackslash}p{#1}}

\makeatletter
\DTMnewdatestyle{dotted}{%
  \renewcommand{\DTMdisplaydate}[4]{\two@digits{##3}.\two@digits{##2}.##1}%
  \renewcommand{\DTMDisplaydate}{\DTMdisplaydate}%
}
\makeatother
\DTMsetdatestyle{dotted}

% redefine spaces after titles
\makeatletter
\renewcommand{\@maketitle}{%
  \begin{center}
    {\huge \bfseries \@title \par}%
    \vskip 0.2em % space between title and author
    {\large \@author \par}%
    % \vskip 0.2em % space between author and date
    % {\normalsize \@date \par}%
    {\small \hfill \begin{english}\DTMtoday\end{english} \par}%
  \end{center}
}
\makeatother

\fancyhf{} % clear default
\fancypagestyle{plain}{
  \fancyhf{}
  \fancyhead[L]{مدرسة التسامح الشاملة}
  % \fancyhead[L]{\includegraphics[height=1cm]{../../../images/logoTasamoh.png}}
  \fancyhead[R]{الأستاذ محمود اغبارية}
  \fancyfoot[C]{\thepage}
}

\fancyhead[L]{مدرسة التسامح الشاملة}
\fancyhead[R]{الأستاذ محمود اغبارية}
\fancyfoot[C]{\thepage}

\setcounter{tocdepth}{3} % only section subsection and subsubsection in TOC

\newcommand{\importpdfpage}[2]{%
    \thispagestyle{empty}
    \newgeometry{left=0cm,right=0cm,top=0cm,bottom=0cm}
    \includegraphics[page=#2,width=0.95\paperwidth,height=0.95\paperheight,keepaspectratio]{#1}
    \restoregeometry
}

\newcommand{\insertFullImg}[1]{%
    \noindent
    \makebox[\textwidth][c]{
        \includegraphics[width=0.9\paperwidth,keepaspectratio]{#1}%
    }%
}%



\title{أكاديمية فرسان التنانين — المرحلة الرابعة: الفئات المجردة والواجهات}

\begin{document}

\maketitle
\thispagestyle{fancy}

\begin{boxWarning}
\textbf{تعليمات هامة للعمل:}
تعتمد هذه المرحلة على الفئات التي بنيتموها في المراحل السابقة (\textenglish{BasicDragon} و \textenglish{FireDragon} و \textenglish{HibernatingDragon} و \textenglish{AlphaDragon} و \textenglish{DragonRegistry}) — حدّثوا الفئات الموجودة ولا تبدأوا من الصفر. \\
في \textbf{كل تمرين}، يجب تحديث مقطع الفحص داخل العملية الرئيسية \textenglish{Main} وتشغيله للتأكد من صحة عمل الفئات والعمليات الجديدة التي كتبتموها.
\end{boxWarning}

% ==============================================================
\section*{المرحلة الرابعة: الفئات المجردة والواجهات}
% ==============================================================

\begin{enumerate}[label=\textbf{تمرين \arabic*:}, wide=0pt, leftmargin=*, itemsep=2.2em]

    \item \textbf{تحويل \textenglish{BasicDragon} إلى فئة مجردة (\textenglish{abstract})} \\
    غيّروا الفئة \textenglish{BasicDragon} بحيث تصبح \textbf{فئة مجردة} (\textenglish{abstract class})، وأضيفوا إليها عملية مجردة تعريفها:
    \begin{boxCode}
    \begin{english}
    \begin{minted}{csharp}
public abstract bool UseEnergy(object o);
    \end{minted}
    \end{english}
    \end{boxCode}
    وغيّروا \textenglish{IsInDanger()} لتصبح عملية مجردة (\textenglish{abstract}) أيضاً. طبّقوا \textenglish{UseEnergy(object o)} في الفئات الوارثة:
    \begin{itemize}
        \item \textenglish{FireDragon} و \textenglish{AlphaDragon}: تتلقى \textenglish{double} (أو \textenglish{int}) وتستهلكه كما في \textenglish{Consume}.
        \item \textenglish{HibernatingDragon}: تتلقى \textenglish{Date} (تاريخ اليوم) وتوقظ التنين كما في \textenglish{WakeUp(today, false)}.
        \item إذا وصل مُعطى من نمط غير مناسب (أو \textenglish{null}) تعيد \textenglish{false}.
    \end{itemize}
    \textbf{الفحص في \textenglish{Main}:} \\
    امسحوا تنانين السجل واستدعوا \textenglish{UseEnergy(25.0)} و \textenglish{UseEnergy(new Date(...))} واطبعوا النتائج.

    \item \textbf{تطبيق الواجهة \textenglish{IComparable}} \\
    طبّقوا الواجهة \textenglish{IComparable} في الفئة \textenglish{Date}، وكذلك في الفئتين:
    \begin{itemize}
        \item \textbf{\textenglish{HibernatingDragon}:} تتم المقارنة حسب \textbf{تاريخ الاستيقاظ} (\textenglish{wakeDate}).
        \item \textbf{\textenglish{AlphaDragon}:} تتم المقارنة حسب \textbf{القوة الكامنة} ($\text{energy} + \text{emberReserve}$).
    \end{itemize}
    \textbf{الفحص في \textenglish{Main}:} \\
    قارنوا بين كائنات من كل فئة باستخدام \textenglish{CompareTo} واطبعوا النتائج.

    \item \textbf{فئة خدمات الفرز (\textenglish{SortServices})} \\
    ابنوا فئة باسم \textenglish{SortServices} وطبّقوا فيها عمليتي فرز ثابتتين لمصفوفة من النمط \textenglish{IComparable[]}:

    {\centering
    \renewcommand{\arraystretch}{1.2}
    \small
    \begin{tabular}{|>{\raggedleft\arraybackslash}p{0.4\textwidth}|>{\raggedright\arraybackslash}p{0.5\textwidth}|}
    \hline
    \multicolumn{2}{|c|}{\textbf{الفئة: \textenglish{SortServices}}} \\
    \hline
    عملية فرز ثابتة ترتّب المصفوفة بترتيب \textbf{تصاعدي}. & \textenglish{public static void SortAscending(IComparable[] arr)} \\
    \hline
    عملية فرز ثابتة ترتّب المصفوفة بترتيب \textbf{تنازلي}. & \textenglish{public static void SortDescending(IComparable[] arr)} \\
    \hline
    \end{tabular}
    \par} %

    \vspace{2em}
    \noindent \textbf{الفحص في \textenglish{Main}:} \\
    افرِزوا مصفوفات من \textenglish{Date[]} و \textenglish{HibernatingDragon[]} و \textenglish{AlphaDragon[]} واطبعوها قبل الفرز وبعده.

    \item \textbf{ترتيب القادة والتنانين النائمة في \textenglish{DragonRegistry}} \\
    أضيفوا العمليتين التاليتين إلى \textenglish{DragonRegistry}:
    \begin{itemize}
        \item \textenglish{GetAlphasByPower()}: تعيد مصفوفة (\textenglish{AlphaDragon[]}) من قادة القطعان مرتّبة حسب \textbf{القوة الكامنة} — الأقوى أولاً (تنازلياً).
        \item \textenglish{GetHibernatingByWakeDate()}: تعيد مصفوفة (\textenglish{HibernatingDragon[]}) من التنانين النائمة مرتّبة حسب \textbf{تاريخ الاستيقاظ} — الأقرب موعداً أولاً (تصاعدياً).
    \end{itemize}
    \textbf{الفحص في \textenglish{Main}:} \\
    استدعوا العمليتين واطبعوا المصفوفتين المرتبتين.

    \item \textbf{تمرين ختامي: محاكاة يوم كامل في الوادي} \\
    اكتبوا في \textenglish{Main} محاكاة ليوم كامل في الأكاديمية تجمع كل ما بنيتموه:
    \begin{enumerate}[label=\textbf{\arabic*.}]
        \item أنشئوا سجلّين (\textbf{سجل الشمال} و\textbf{سجل الجنوب}) في كل منهما تنانين من الأنماط الثلاثة ويشتركان في تنينين على الأقل، ثم ادمجوهما في \textbf{سجل الوادي} عبر \textenglish{Merge}.
        \item استدعوا \textenglish{FeedAll} بكمية صغيرة واطبعوا \textenglish{TotalEnergy()}.
        \item أرسلوا ثلاثة تنانين نارية في رحلات طيران (\textenglish{FlyTo}) بحيث يدخل اثنان منهم في جمر الطوارئ، ثم اطبعوا \textenglish{DangerLevel()}، \textenglish{DangerDragons()}، و \textenglish{PrintSuggestions()}.
        \item أجروا إنقاذاً عبر \textenglish{TransferEnergy} من قائد القطيع إلى التنانين النارية التي في خطر، ثم أعلنوا إنذار بركان وأيقظوا التنانين النائمة عبر \textenglish{WakeUp(today, true)}.
        \item اطبعوا قادة القطعان مرتبين بالقوة الكامنة، والتنانين النائمة حسب موعد الاستيقاظ، ثم احذفوا التنانين التي استيقظت واطبعوا السجل النهائي.
    \end{enumerate}

\end{enumerate}

\section*{خريطة المراجعة للمعلم والطالب}
\begin{center}
\renewcommand{\arraystretch}{1.35}
{\small
\begin{tabular}{|p{0.58\textwidth}|p{0.34\textwidth}|}
\hline
\textbf{الموضوع المُراجَع} & \textbf{أرقام التمارين} \\
\hline
عمليات مساعدة (\textenglish{private}) لتجنب تكرار الكود & 1، 3، 7، 8، 16 \\
\hline
المصفوفات (إنشاء، مسح، فحص \textenglish{null}، بحث، إزاحة، دمج) & 3، 7، 8، 14، 16، 20 \\
\hline
فئة تحوي كائناً مركباً أو مصفوفة كائنات & 2، 3، 4، 7، 8 \\
\hline
الشروط المركّبة (\textenglish{\&\&}, \textenglish{||}, \textenglish{!}) والشروط المتداخلة & 1، 2، 4، 7، 9، 12 \\
\hline
العمليات البنائية الناسخة والمؤشرات & 2، 3، 4، 16 \\
\hline
الوراثة، \textenglish{base}، وأذونات الوصول \textenglish{protected} & 6، 7، 10 \\
\hline
إعادة التعريف (\textenglish{virtual}/\textenglish{override}) وفحص الأنماط (\textenglish{is}/\textenglish{as}) & 9–13 \\
\hline
العملية \textenglish{Equals}، الفئات المجردة (\textenglish{abstract})، والواجهة \textenglish{IComparable} & 15، 17–20 \\
\hline
\end{tabular}
}
\end{center}

\vspace{1cm}
\begin{flushleft}
أرجو لكم وقتًا ممتعًا.

الأستاذ محمود اغبارية.
\end{flushleft}

\end{document}
